\section{CpuLocal、原生系统调用与安全返回}

v0.8 第一次进入 Ring 3 时，项目选择 \code{INT 0x80}。这是有意的学习顺序：
异常章已经建立 GDT、IDT、TSS 和 \code{iretq}，软件中断只需把一扇 IDT gate
的 DPL 设为 3，就能复用同一套特权切换与保存帧。CPU 会自动读取 TSS.RSP0，
从用户栈切换到当前 Thread 的 Ring 0 栈，再压入
\code{SS:RSP:RFLAGS:CS:RIP}。入口安全性大部分由已经理解的硬件机制提供。

原生 \code{syscall/sysret} 的设计目标不同。AMD 在 x86-64 中把频繁特权转换
压缩成少量寄存器和 MSR 操作，避免通用 IDT gate 的描述符遍历。Intel 后来
也在 64 位模式采用这组指令。硬件做得更少，操作系统必须负责更多：

\begin{itemize}
  \item \code{syscall} 不查询 TSS，也不修改 RSP；
  \item 用户 RIP 与 RFLAGS 分别进入 RCX、R11，而不是自动压栈；
  \item CS/SS 与入口 RIP 来自 STAR/LSTAR；
  \item 活动 RFLAGS 按 FMASK 清位；
  \item 用户与内核的 GS base 需要软件交换；
  \item \code{sysret} 前必须证明 RCX、R11 和段状态安全。
\end{itemize}

\begin{keyidea}
\code{syscall} 不是“更快的 \code{int 0x80}”。它把换栈、每 CPU 状态、统一
保存现场和安全返回从硬件门描述符移交给内核。只有 Thread 栈所有权已经冻结，
才能单独验收这份新合同。
\end{keyidea}

\topic{三代跨特权入口背后的取舍}
保护模式提供 call gate。GDT/LDT 描述符可以指定目标段、偏移、目标特权级和
需要复制的参数数量，功能通用但规则繁复。32 位系统又使用
\code{SYSENTER/SYSEXIT}，它通过专用 MSR 缩短路径，却没有天然适配长模式的
完整返回语义。软件中断则最容易调试，但需要 IDT gate 和通用中断路径。

\begin{osgridlongtable}{L{0.18\textwidth}L{0.23\textwidth}L{0.24\textwidth}L{0.25\textwidth}}
机制 & 入口来源 & 是否自动换栈 & 返回与主要代价\\ \hline
call gate & GDT/LDT gate & 按 TSS 与 gate 规则 & RETF；描述符语义复杂\\ \hline
\code{INT n} & IDT interrupt/trap gate & 是，使用 TSS.RSP0 & IRETQ；通用且路径较长\\ \hline
SYSENTER & SYSENTER MSR & 不自动保存完整返回现场 & SYSEXIT；主要面向 32 位\\ \hline
SYSCALL & STAR/LSTAR/FMASK & 否 & SYSRETQ；软件承担可信栈与返回校验\\ \hline
\end{osgridlongtable}

本项目保留 \code{INT 0x80}，不是因为不敢删除旧路径，而是把它变成差分基准。
同一个 PID 查询分别经两条入口，若返回值或副作用不同，错误一定发生在进入/
规范化/返回边界，而不是系统调用业务本身。

\topic{先冻结处理器规格，再执行依赖能力的指令}
“目标是 x86-64”仍不足以说明一台 CPU 暴露哪些可选能力。长模式、NX、
FXSAVE、SSE2 和 SYSCALL 分布在不同 CPUID 叶；地址宽度又决定规范地址和页表
能表达的范围。若代码先执行 \code{wrmsr} 或 \code{syscall}，再等待 \#UD
暴露能力缺失，系统已经进入难以解释的半初始化状态。

v1.3 读取三组叶：

\begin{osgridlongtable}{L{0.24\textwidth}L{0.16\textwidth}L{0.21\textwidth}L{0.27\textwidth}}
CPUID 位置 & 位 & 名称 & 本阶段用途\\ \hline
\code{01H:EDX} & 24 & FXSR & FXSAVE/FXRSTOR 存在\\ \hline
\code{01H:EDX} & 25 & SSE & XMM/MXCSR 架构状态存在\\ \hline
\code{01H:EDX} & 26 & SSE2 & 当前 64 位用户指令基线\\ \hline
\code{80000001H:EDX} & 11 & SYSCALL & 原生入口/返回存在\\ \hline
\code{80000001H:EDX} & 20 & NX & 页表执行禁止能力存在\\ \hline
\code{80000001H:EDX} & 29 & LM & long mode 存在\\ \hline
\code{80000008H:EAX} & 7:0 & physical width & 必须位于 36--52\\ \hline
\code{80000008H:EAX} & 15:8 & linear width & 当前精确要求 48\\ \hline
\end{osgridlongtable}

六项布尔能力被编码为 \code{uint64_t} mask：

\[
  R = 2^0+2^1+2^2+2^3+2^4+2^5 = \code{0x3F}.
\]

\code{ProcessorFeatureProfile} 同时保存 available、missing、两个宽度和逐项
布尔值。missing 由
\[
  M = R \mathbin{\&} \neg A
\]
一次计算得到。日志既给出可读的阶段状态，也给出机器可解析的位图；失败用例
禁用 SYSCALL 时，\(M=\code{0x20}\)。

\begin{contractbox}
最大 CPUID 叶必须先检查。若 CPU 没有声明 \code{80000008H}，宽度不能用零
以外的猜测值替代。零宽度进入明确的 profile 失败，不继续使用宿主或 QEMU
型号的经验默认值。
\end{contractbox}

\topic{为什么当前只接受 48 位，不顺便开放 LA57}
四级页表使用虚拟地址 47--12 位形成 PML4/PDPT/PD/PT 四个索引，bit 47 以上
必须是符号扩展。五级页表把实现宽度扩展到 57 位，但需要在进入分页前设置
CR4.LA57，内核/用户地址布局、canonical 检查和页表管理都要共同变化。

仅因为 CPUID 报告 LA57 就把宽度传入一个通用函数，会让软件宣称自己支持从未
建立的第五级。因此 profile 当前把 linear width 精确冻结为 48。未来开放
57 位必须作为独立阶段，改变启动链、内存布局、测试向量和文档，而不是放宽一个
比较符。

\topic{六个 MSR 共同定义原生入口}
CPUID 只证明能力存在。操作系统还要通过 \code{RDMSR/WRMSR} 配置：

\begin{osgridlongtable}{L{0.22\textwidth}L{0.17\textwidth}L{0.49\textwidth}}
MSR & 地址 & 当前合同\\ \hline
IA32\_EFER & \code{0xC0000080} & 保留 LME/LMA/NXE，设置 SCE bit 0\\ \hline
IA32\_STAR & \code{0xC0000081} & kernel selector base 与 SYSRET selector base\\ \hline
IA32\_LSTAR & \code{0xC0000082} & 64 位入口 \code{OsKernelNativeSystemCallEntry}\\ \hline
IA32\_FMASK & \code{0xC0000084} & 进入时清 TF、IF、DF、NT、AC\\ \hline
IA32\_GS\_BASE & \code{0xC0000101} & 当前用户 GS base 为 0\\ \hline
\shortstack[l]{IA32\_KERNEL\_\\GS\_BASE} &
\code{0xC0000102} & \code{CpuLocal} 地址\\ \hline
\end{osgridlongtable}

\code{native_system_call_layout.cpp} 是不执行特权指令的纯布局层。它验证地址
宽度、selector 顺序与两个规范地址，再构造期望值；目标机层负责写 MSR，
随后把六项全部读回并逐项比较。这样“selector 算错”和“硬件写回不符”分别
可由宿主集成测试与 QEMU 定位。

\topic{STAR 的高字段不是最终用户 CS}
当前 GDT 顺序为：

\begin{lstlisting}[caption={与 SYSCALL/SYSRET 配套的 GDT 选择子}]
kernel code = 0x08
kernel data = 0x10
user data   = 0x18 | 3 = 0x1B
user code   = 0x20 | 3 = 0x23
\end{lstlisting}

SYSCALL 使用 STAR[47:32] 作为内核 CS，SS 隐式取 CS+8。64 位 SYSRET 对
STAR[63:48] 的规则不同：

\[
\begin{aligned}
\mathrm{SS}_{user} &= (\mathrm{base}+8)\mathbin{|}3,\\
\mathrm{CS}_{user} &= (\mathrm{base}+16)\mathbin{|}3.
\end{aligned}
\]

令 base=\code{0x10}，才得到 SS=\code{0x1B}、CS=\code{0x23}。于是：

\[
  \mathrm{STAR}
  = (\code{0x10}\ll48)\mathbin{|}(\code{0x08}\ll32)
  = \code{0x0010000800000000}.
\]

\begin{failurebox}
把最终用户代码选择子 \code{0x23} 原样写入 STAR 高 16 位，会让 SYSRET 再
加 16，选择完全不同的 GDT 项。日志看起来只是一个十六进制值，错误却要到
第一次返回才以 \#GP 或三重故障表现，所以必须在写 MSR 前验证选择子算术。
\end{failurebox}

\topic{FMASK 改变内核活动标志，不抹掉用户原值}
SYSCALL 先把用户 RFLAGS 复制到 R11，再从当前活动 RFLAGS 清除 FMASK 指定
的位。当前 \code{FMASK=0x44700}：

\begin{osgridlongtable}{L{0.12\textwidth}L{0.18\textwidth}L{0.58\textwidth}}
位 & 名称 & 入口必须清除的原因\\ \hline
8 & TF & 防止用户单步在可信栈形成前触发调试异常\\ \hline
9 & IF & 完整保存现场前不能接受可屏蔽 IRQ\\ \hline
10 & DF & C++ 与字符串操作要求向高地址方向\\ \hline
14 & NT & 不继承遗留嵌套任务语义\\ \hline
18 & AC & 不把用户对齐检查状态带入内核\\ \hline
\end{osgridlongtable}

R11 中仍是用户原值。入口 \code{cld} 只维护当前内核 ABI，构造
\code{UserContext} 时写入的是 R11。因此设置 DF 的测试能证明“内核执行时
清 DF”和“返回时保留用户 DF”不是相互冲突的要求。

\topic{CpuLocal 把每 CPU 入口状态收束为一个对象}
TSS.RSP0 属于 CPU，不属于 Process；原生入口栈同样是“当前 CPU 正在运行
哪个 Thread”的派生状态。若 current Thread、TSS.RSP0、原生 entry RSP 和
need-resched 分散在无关全局变量中，调度切换很难证明它们一起更新。

v1.3 的 \code{CpuLocal} 按缓存线大小 64 字节对齐。前四个字段是 C++ 与汇编
共同遵守的 ABI：

\begin{osgridlongtable}{L{0.13\textwidth}L{0.31\textwidth}L{0.44\textwidth}}
偏移 & 字段 & 不变量\\ \hline
0 & self address & 初始化后等于对象自身地址\\ \hline
8 & current Thread index & 与 scheduler current 一致；无 Thread 为 UINT64\_MAX\\ \hline
16 & kernel entry RSP & 当前 Thread 动态栈顶，非零且 16 B 对齐\\ \hline
24 & syscall user RSP scratch & 只由原生入口换栈前暂存，退出入口时清零\\ \hline
\end{osgridlongtable}

后续字段保存 IRQ/抢占深度与峰值、need-resched、活动系统调用来源，以及入口、
IRQ 打断、返回前调度、SYSRET/IRET、拒绝和可信栈校验统计。
\code{static_assert(offsetof(...))} 把 C++ 布局锁定；NASM 仍用全大写
\code{equ} 常量表达不可由 C++ 类型系统共享的汇编偏移。

\begin{pdfdiagram}
  \node[os hardware] (cpu) {当前 BSP\\GS base};
  \node[os state,right=of cpu] (local) {CpuLocal\\current / entry RSP\\depth / need};
  \node[os block,above right=of local] (tss) {TSS.RSP0\\INT/IRQ 自动换栈};
  \node[os block,below right=of local] (syscall) {GS:entry RSP\\SYSCALL 手动换栈};
  \node[os state,right=of tss] (thread) {Thread 动态栈\\4 页 RW/NX\\双 guard};
  \draw[os arrow] (cpu) -- node[above]{SWAPGS} (local);
  \draw[os arrow] (local) -- (tss);
  \draw[os arrow] (local) -- (syscall);
  \draw[os arrow] (tss) -- (thread);
  \draw[os arrow] (syscall) -- (thread);
\end{pdfdiagram}
\diagramnote{两种入口使用同一 Thread 栈；TSS 与 CpuLocal 是两个硬件入口的适配状态，不是两份独立所有权。}

\topic{深度为什么比布尔值更可靠}
IRQ 可以嵌套，抢占禁止 guard 也可能由调用链分层取得。单个布尔值不能区分
“内层释放”与“最后一层释放”。深度从零进入一、离开时减一；零下溢和
\code{UINT64\_MAX} 上溢都明确失败。最大深度只在新峰值出现时更新，运行结束
要求当前深度归零但峰值非零。

\code{CpuPreemptionGuard} 构造时增加深度，析构时减少；失败不伪造取得成功。
它没有让内核自动变成 SMP-safe，只把当前单 BSP 的合法调度点变成可检查状态。

\topic{原生入口逐条建立可信状态}
用户包装器按 System V ABI 收到 number 与四个参数，再变换为项目 ABI：

\begin{lstlisting}[language={[x86masm]Assembler},caption={用户侧最小原生包装}]
mov rax, rdi        ; number
mov rdi, rsi        ; argument 0
mov rsi, rdx        ; argument 1
mov rdx, rcx        ; argument 2
mov r10, r8         ; argument 3
syscall
ret
\end{lstlisting}

SYSCALL 刚到 LSTAR 时，CPL 已是 0，但 RSP 仍完全由用户控制。入口不能
\code{push}，也不能调用任何函数。当前顺序为：

\begin{lstlisting}[language={[x86masm]Assembler},caption={可信换栈的关键前缀}]
swapgs
mov [gs:24], rsp
mov rsp, [gs:16]

push 0x1B          ; user SS
push qword [gs:24] ; saved user RSP
push r11           ; saved user RFLAGS
push 0x23          ; user CS
push rcx           ; saved user RIP
push 0             ; normalized error code
push 0x81          ; native entry vector
\end{lstlisting}

只有切到可信 RSP 后才压栈。随后按与通用中断相同的顺序保存 15 个 64 位通用
寄存器。保存完成后才能加载内核数据段、按 16 字节对齐 C++ 调用栈，并
\code{sti} 开放 IRQ。

\begin{contractbox}
GS:24 是暂存，不是用户指针可信化。C++ 仍要验证它等于 UserContext 中的
RSP，并检查该 RSP 属于当前用户栈映射。暂存只能证明“汇编没有在换栈途中取错
值”，不能证明用户地址本身安全。
\end{contractbox}

\topic{SWAPGS 为什么能在不占通用寄存器时找到 CpuLocal}
用户态时 IA32\_GS\_BASE 当前为零，IA32\_KERNEL\_GS\_BASE 保存 CpuLocal
地址。SWAPGS 原子交换二者，因而第一条后 GS 就能寻址内核对象。若先把一个
通用寄存器用于地址，必须先保存它；而保存又需要可信内存，形成循环依赖。

未来用户 GS/TLS 非零时，逻辑不变：用户 base 被暂存在 KERNEL\_GS\_BASE
一侧，返回前第二次 SWAPGS 恢复。当前不开放用户 GS，是明确阶段边界，不是
SWAPGS 只能处理零值。

\topic{两条入口规范化为同一份 UserContext}
原 \code{ExceptionFrame} 保存 15 个通用寄存器、vector、error code、
RIP、CS 与 RFLAGS，共 160 字节。来自 Ring 3 的现场还具有 RSP 与 SS，
统一结构因此为 176 字节：

\begin{osgridlongtable}{L{0.17\textwidth}L{0.17\textwidth}L{0.26\textwidth}L{0.28\textwidth}}
偏移 & 长度 & 字段组 & 来源\\ \hline
0--119 & 120 B & R15..RAX & 汇编逐项保存/人工初始现场\\ \hline
120--127 & 8 B & vector & Initial、IRQ、INT80、SYSCALL\\ \hline
128--135 & 8 B & normalized error & 系统调用与 IRQ 为零\\ \hline
136--159 & 24 B & RIP、CS、RFLAGS & CPU 帧或 RCX/R11 规范化\\ \hline
160--175 & 16 B & RSP、SS & CPU 特权帧或原生入口人工补齐\\ \hline
\end{osgridlongtable}

vector 不表示“发生异常”，而是统一现场的来源标签：

\begin{osgridlongtable}{L{0.28\textwidth}L{0.20\textwidth}L{0.39\textwidth}}
来源 & vector & 最终返回基线\\ \hline
首次进入用户态 & 0 & IRETQ\\ \hline
8259A IRQ & 32--47 & IRETQ\\ \hline
\code{INT 0x80} & \code{0x80} & IRETQ\\ \hline
\code{SYSCALL} & \code{0x81} & 满足快速白名单时 SYSRETQ，否则 IRETQ\\ \hline
\end{osgridlongtable}

系统调用 C++ switch 只存在一份。入口差异在进入分发器前消失；阻塞、退出、
用户复制、文件和管道逻辑不需要知道最初执行了哪条指令。来源只服务计数、入口
一致性和返回选择。

\topic{返回前的校验顺序是一条安全证明}
SYSRET 的 RCX 来自用户 RIP。非规范地址可能在处理器仍处于 Ring 0 的阶段
触发 \#GP，因此不能把“让硬件试一下”当作验证。当前返回准备从所有权开始：

\begin{enumerate}[label=\arabic*.]
  \item 调度必须活动，frame 必须属于 CpuLocal.current Thread；
  \item frame 的 176 字节必须完整位于该 Thread 活动动态内核栈；
  \item RIP 与 RSP 必须是 48 位低半规范地址；
  \item RIP 查询当前 CR3 后必须是 user、executable、not writable；
  \item 栈探测地址必须是 user、writable、not executable；
  \item CS 必须为 \code{0x23}，SS 必须为 \code{0x1B}；
  \item RFLAGS bit 1 必须为 1，未开放或危险位必须为零。
\end{enumerate}

\topic{规范地址与低半用户地址不是同一概念}
48 位模式以 bit 47 作为符号位：

\[
\begin{aligned}
\text{lower} &= [0,\code{0x00007FFFFFFFFFFF}],\\
\text{hole}  &= [\code{0x0000800000000000},
                  \code{0xFFFF7FFFFFFFFFFF}],\\
\text{upper} &= [\code{0xFFFF800000000000},2^{64}-1].
\end{aligned}
\]

lower 与 upper 都是 canonical，hole 不是。用户返回只允许 lower；否则一个
完全规范的高半内核地址仍可能被当作用户 RIP。通过低半检查后还要查 PTE，
因为“地址形状合法”不表示页面存在或具有用户权限。

\topic{RFLAGS 使用合法集合与快速集合}\par
bit 1 是架构保留常一位。CF/PF/AF/ZF/SF/TF/IF/DF/OF/RF 等被当前模型明确
分类；IOPL、NT、VM、AC、ID 等不开放。合法并不等于适合 SYSRET：

\begin{itemize}
  \item 普通原生现场可走 SYSRET；
  \item DF 表示用户方向状态，当前用 IRET 精确恢复；
  \item RF 可能由硬件调试/异常返回语义带入，也用 IRET；
  \item INT80/IRQ/Initial 即使标志简单，仍使用它们原本的 IRET 返回路径。
\end{itemize}

\begin{pdfdiagram}
  \node[os state] (frame) {候选 UserContext};
  \node[os block,right=of frame] (validate) {所有权 + canonical\\映射 + 段 + RFLAGS};
  \node[os hardware,above right=of validate] (sysret) {Native + fast\\SYSRETQ};
  \node[os hardware,right=of validate] (iret) {合法非快速\\IRETQ};
  \node[os state,below right=of validate] (reject) {非法\\终止用户进程};
  \draw[os arrow] (frame) -- (validate);
  \draw[os arrow] (validate) -- (sysret);
  \draw[os arrow] (validate) -- (iret);
  \draw[os arrow] (validate) -- (reject);
\end{pdfdiagram}
\diagramnote{快速路径是合法集合的真子集；回退 IRET 是正常语义，不是失败。}

\topic{非法返回为什么终止用户进程而不是内核 panic}
如果 frame 确实属于当前用户 Thread，但字段被破坏，错误的安全域是该用户
执行流。内核输出一次有界诊断，记录 ownership、validation status、mapping、
entry、vector、RIP、RSP 与 RFLAGS，再按用户异常终止该 Process。调度器选择
下一份 frame 后，返回准备循环重新验证。

只有 frame 指针为空、CpuLocal 状态损坏、计数溢出等内核自身不变量失败才
停机。把用户输入错误和内核内部错误分开，才能避免一个普通用户现场导致整个
系统 panic。

\topic{IRQ 可以打断系统调用，但不能在任意内核调用链换栈}
完整 UserContext 形成后，入口执行 \code{sti}，因此 1000 Hz PIT 可以在
系统调用 C++ 分发期间到达。IRQ 进入/离开更新 CpuLocal 深度；若活动
system-call depth 为一，再累计 \code{interrupt_during_system_call_count}。

Ring 0 中的 timer IRQ 不直接调用完整换栈调度。它只设置
\code{need_reschedule=1}，EOI 后回到原系统调用。分发结束、入口执行
\code{cli} 后，\code{OsKernelPrepareUserReturn} 消费请求，在这个已知边界
调用 scheduler：

\begin{pdfdiagram}
  \node[os block] (dispatch) {syscall C++\\IF=1};
  \node[os hardware,right=of dispatch] (irq) {PIT IRQ0\\depth++};
  \node[os state,right=of irq] (request) {need-resched=1\\EOI, depth--};
  \node[os block,below=of request] (prepare) {return prepare\\IF=0};
  \node[os state,left=of prepare] (schedule) {Thread scheduler\\选择 resume frame};
  \node[os hardware,left=of schedule] (return) {validate\\SYSRET/IRET};
  \draw[os arrow] (dispatch) -- (irq);
  \draw[os arrow] (irq) -- (request);
  \draw[os arrow] (request) -- (prepare);
  \draw[os arrow] (prepare) -- (schedule);
  \draw[os arrow] (schedule) -- (return);
\end{pdfdiagram}
\diagramnote{IRQ 提交事实，返回边界执行策略；任何普通内核函数都不会突然失去自己的调用栈。}

为让这条证据不是偶然时序，第一个日志系统调用在完整入口内等待 timer tick
变化，并使用 \code{hlt} 让真实 IRQ 唤醒。QEMU 必须看到
\code{SYSCALL_IRQ_INTERRUPTS>0}。IRQ 又必须设置调度请求，最终
\code{SYSCALL_RETURN_RESCHEDULES>0}。两项分别证明“中断发生”和“请求在正确
边界消费”。

\topic{非局部返回暴露了 SWAPGS 的真正状态所有权}
正常原生尾部在恢复用户段之前执行第二次 SWAPGS。可是系统调用可能使最后一个
可运行 Thread 阻塞或退出；调度器没有后继时，会直接恢复
\code{ExecuteProcesses} 之前保存的永久内核调用栈，而不是回到系统调用汇编
尾部。

若仍假设“每次 SYSCALL 都自然经过一次 SYSRET”，GS 会停留在内核一侧。
下一次用户原生入口执行 SWAPGS，反而把零用户 base 换入活动 GS，随后
\code{[gs:16]} 访问错误地址。表面症状可能是第二轮 Shell 输入后无日志，
根因却发生在上一次 Thread 非局部退出。

修复顺序是：

\begin{enumerate}[label=\arabic*.]
  \item 在清理 CpuLocal 前读取 \code{NativeSystemCallActive()}；
  \item 把结果编码为显式 64 位 C ABI 参数；
  \item 恢复永久栈的汇编桩先按参数补做 SWAPGS；
  \item 再恢复内核数据段和保存的寄存器调用链；
  \item \code{ClearCurrentThread} 清入口 depth、来源和用户 RSP scratch。
\end{enumerate}

不能从最终 resume frame 的 vector 推断是否需要 SWAPGS：返回前调度可能已经
选择另一 Thread 的 IRQ 或初始 frame，但当前硬件 GS 状态仍由“从哪条入口进入
这个内核调用链”决定。

\topic{NMI 必须承认自己可能落在最危险的两条指令之间}
IF=0 不能屏蔽 NMI。NMI 可能发生在 SYSCALL 刚进入、尚未 SWAPGS，也可能发生
在已经 SWAPGS、尚未切到可信 RSP。普通 C++ handler 若无条件读 GS 或使用当前
栈，至少一种状态会错。

成熟内核可实现 paranoid entry：根据保存 CS、GS base 或嵌套标记判断交换
状态，并使用独立 IST 栈完成可恢复处理。该协议本身足以构成一个独立阶段。
v1.3 选择较小但真实的边界：

\begin{itemize}
  \item IDT vector 2 使用 TSS.IST2；
  \item 汇编桩不访问普通 RSP/GS/CpuLocal；
  \item 只把一个固定 BSS 观察位写为一；
  \item 关闭中断并永久 HLT；
  \item 不打印、不分配、不取锁、不尝试返回。
\end{itemize}

\begin{failurebox}
“NMI 很少发生”不能证明入口正确。当前 fail-stop 不提供恢复能力，但明确避免
在未知 GS/栈状态中进入复杂运行时。以后实现可恢复 NMI 时，必须先替换这份
合同，而不是直接在最小桩中调用普通异常分发器。
\end{failurebox}

\topic{用户侧三个标记分别回答三个问题}
worker 执行：

\begin{enumerate}[label=\arabic*.]
  \item 用原生和兼容入口各查询一次 PID，比较返回值，输出
    \code{DUAL_SYSCALL_ENTRY_EQUIVALENT}；
  \item 用普通原生入口完成一次调用并继续执行，输出
    \code{SYSRET_RETURNED}；
  \item 在原生入口前 \code{std}，返回后 \code{cld}，确认结果并输出
    \code{SYSCALL_IRET_FALLBACK_RETURNED}。
\end{enumerate}

第二个标记本身不能直接观察 CPU 用了 SYSRET 还是 IRET，所以必须与 Kernel
\code{SYSRET_RETURNS>0} 组合。第三个标记与
\code{SYSCALL_IRET_FALLBACKS=1} 组合，证明 DF 不是在分发器中被永久丢弃。
双入口调用选无副作用 PID 查询，避免“调用两次本来就应该改变状态”使差分失去
意义。

\topic{四层测试从不同位置观察同一合同}\par
\begin{osgridlongtable}{L{0.17\textwidth}L{0.30\textwidth}L{0.41\textwidth}}
层次 & 主要输入 & 能证明的事实\\ \hline
单元 & CPUID 叶、UserContext 边界 & profile、canonical、段、RFLAGS 和返回选择\\ \hline
集成 & CpuLocal 与六 MSR 纯布局 & 深度/统计/清理、STAR 算术与回读比较\\ \hline
随机 & 100000 个确定性现场 & 400000 项规范性、合法性、选择和拒绝性质\\ \hline
QEMU & 真实 SYSCALL、SYSRET、IRQ 与 MSR &
可信换栈、双入口、IRQ 打断、调度与失败停止\\ \hline
\end{osgridlongtable}

随机种子为 \code{0x5A17C011BADC0FFE}。每轮随机地址先与独立 48 位 canonical
定义比较，再构造一个合法 Initial/IRQ/INT80/SYSCALL 现场；返回选择必须与
独立白名单条件一致。最后强制把 RIP 置入地址空洞，验证拒绝。每轮四项性质，
合计 400000 项。

\topic{QEMU 冷路径统计如何判读}
一次代表性 256 MiB functional 运行得到：

\begin{lstlisting}[caption={v1.3 原生入口代表性统计}]
LEGACY_SYSCALL_ENTRIES=0x2
NATIVE_SYSCALL_ENTRIES=0x230
SYSCALL_IRQ_INTERRUPTS=0x56
SYSCALL_RETURN_RESCHEDULES=0x36
SYSRET_RETURNS=0x1C1
IRET_RETURNS=0xE
SYSCALL_IRET_FALLBACKS=0x1
REJECTED_USER_RETURNS=0x0
TRUSTED_SYSCALL_STACK_VALIDATIONS=0x232
\end{lstlisting}

PIT 与宿主调度会改变累计值，测试不把 \code{0x230} 等写成精确常数；它要求
语义下界。兼容与原生入口、IRQ、reschedule、两类返回必须非零，fallback
至少一，拒绝精确为零，可信栈验证不少于入口总数。结束 IRQ/抢占深度和
need-resched 精确为零。

失败路径用 \code{-cpu qemu64,-syscall} 启动同一 ROM、Stage 1、磁盘和
Kernel payload。它必须输出 missing mask \code{0x20}，且禁止出现 processor
ready、扩展现场、GDT、用户态、panic 和最终 READY。这证明系统主动拒绝能力，
不是执行未支持指令后偶然 \#UD。

\topic{调试 ELF 与启动 payload 分离是本阶段的工程修复}
随着 C++ 类型、测试行号和文档关联增加，Debug DWARF 会让
\filepath{kernel.elf} 文件长度持续增长。Stage 1 的内核载荷区位于文件系统
LBA 2048 之前；若直接把完整 Debug 文件写盘，非加载调试段也会占用磁盘扇区，
最终与文件系统重叠。删除调试信息会让 GDB 失去源码、类型和行号。

当前构建保留两类产物：

\begin{itemize}
  \item \filepath{kernel.elf}：LLD 原始 Debug ELF，用于 GDB、符号和开发审计；
  \item \filepath{kernel.payload.elf}：由
    \code{llvm-objcopy --strip-debug} 生成，写入启动盘。
\end{itemize}

strip 只删除不进入运行时的调试段，三个 \code{PT_LOAD}、入口与目标代码内容
不变。正常、Ring 0 故障和 Ring 3 隔离内核变体使用同一个 CMake 生成函数，
避免每次打包临时扩大分区或手工删除文件。Stage 1 审计实际写盘 payload，
GDB 仍加载完整 ELF。

\topic{按目录走读 v1.3}\par
\begin{osgridlongtable}{L{0.43\textwidth}L{0.45\textwidth}}
文件 & 阅读时要回答的问题\\ \hline
\filepath{arch/processor_features.hpp/.cpp} &
CPUID 最大叶、六位 mask 和地址宽度为何是一份 profile？\\ \hline
\filepath{arch/cpu_local.hpp/.cpp} &
汇编 ABI 字段、current Thread、深度和统计如何保持一致？\\ \hline
\filepath{arch/native_system_call_layout.hpp/.cpp} &
为什么 STAR 用户基值是 \code{0x10}，六项回读如何比较？\\ \hline
\filepath{arch/native_system_call.hpp/.cpp} &
纯布局怎样落实为真实 RDMSR/WRMSR，又如何区分失败状态？\\ \hline
\filepath{arch/user_context.hpp/.cpp} &
四类来源如何共享 176 字节布局，合法集合与快速集合如何区分？\\ \hline
\filepath{arch/architecture.asm} &
SWAPGS、换栈、STI 时刻、返回弹栈和非局部恢复分别在哪里？\\ \hline
\filepath{user/system_calls.cpp} &
入口所有权、映射验证、need-resched、非法返回循环和返回选择如何分层？\\ \hline
\filepath{process/process_runtime.cpp} &
TSS/CpuLocal 如何同步，Thread 切换和永久栈恢复如何收束 GS？\\ \hline
\filepath{user/src/system_call.asm} &
默认、兼容和 DF 测试包装如何保持同一参数 ABI？\\ \hline
\end{osgridlongtable}

推荐先读三个纯策略文件及宿主测试，再读 CpuLocal，最后逐条跟踪汇编与
ProcessRuntime。若直接从 \code{sysret} 向前逆推，会同时遇到 selector、
RFLAGS、调度、页表、栈和 GS 六种状态，容易把结果当成原因。

\topic{本阶段明确没有做什么}
当前 CpuLocal 只有一个 BSP 实例；没有 AP 启动、per-CPU 分配、跨核 IPI 或
远端 TLB shootdown。用户 GS/TLS 尚未开放，IA32\_GS\_BASE 为零。虚拟地址
继续固定四级 48 位，LA57 未启用。NMI 采用 fail-stop，不提供可恢复诊断。
\code{INT 0x80} 继续作为兼容入口，移除时间留给 ABI 冻结阶段。

本章完成时，系统调用背后的描述符仍是 v1.0 固定槽；紧随其后的 v1.4 已经
建立类型化 KernelObject、动态 FileTable 和共享 FileDescription，而且没有
改写本章冻结的架构入口、UserContext 和返回验证。这个后续结果反向证明了
分层的意义：对象层可以演进，特权边界不需要跟着重新冒险。

\begin{keyidea}
v1.3 的完成不是“用户程序执行了一次 \code{syscall}”。完整结论是：能力在
使用前被证明，入口在接触内存前取得可信栈，两条入口汇聚为一份现场，IRQ 只在
合法边界请求调度，返回指令由白名单选择，非局部控制流与 NMI 都有明确 GS/栈
策略，而且正常、随机与能力失败证据可以从复位向量重新产生。
\end{keyidea}
